\documentclass[11pt,letterpaper]{article}
\usepackage[T1]{fontenc}
\usepackage[utf8]{inputenc}
\usepackage{amsmath,amsthm,mathtools}
\usepackage{newtxtext,newtxmath}
\usepackage[margin=1in,headheight=15pt]{geometry}
\usepackage{microtype}
\usepackage{booktabs,array,tabularx,longtable}
\usepackage{xcolor}
\usepackage{enumitem,fancyhdr,titlesec}
\usepackage[most]{tcolorbox}
\usepackage{hyperref}
\usepackage{cleveref}
\definecolor{ink}{HTML}{153247}
\definecolor{accent}{HTML}{176B75}
\definecolor{pale}{HTML}{F0F6F7}
\hypersetup{colorlinks=true,linkcolor=accent,citecolor=accent,urlcolor=accent,
 pdftitle={Elliptic addition and Jacobi continued fractions},
 pdfauthor={Research report prepared with ChatGPT},
 pdfsubject={An explicit proof of Barry's corrected elliptic continued-fraction conjecture}}
\titleformat{\section}{\large\bfseries\color{ink}}{\thesection}{0.65em}{}
\titleformat{\subsection}{\normalsize\bfseries\color{ink}}{\thesubsection}{0.65em}{}
\pagestyle{fancy}
\fancyhf{}
\fancyhead[L]{\small\color{ink}Elliptic addition and Jacobi fractions}
\fancyhead[R]{\small\color{ink}Research report}
\fancyfoot[C]{\small\thepage}
\renewcommand{\headrulewidth}{0.3pt}
\setlength{\parindent}{1.2em}
\setlength{\parskip}{3pt}
\setlength{\emergencystretch}{2em}
\numberwithin{equation}{section}
\newtheorem{theorem}{Theorem}[section]
\newtheorem{proposition}[theorem]{Proposition}
\newtheorem{lemma}[theorem]{Lemma}
\newtheorem{corollary}[theorem]{Corollary}
\theoremstyle{definition}
\newtheorem{definition}[theorem]{Definition}
\newtheorem{example}[theorem]{Example}
\theoremstyle{remark}
\newtheorem{remark}[theorem]{Remark}
\newcommand{\Q}{\mathbb{Q}}
\newcommand{\Z}{\mathbb{Z}}
\newcommand{\K}{K}
\newcommand{\OO}{\mathcal{O}}
\newcommand{\D}{\mathcal{D}}
\newcommand{\coef}[2]{[z^{#1}]\,#2}
\newcommand{\inv}{\mathcal{I}}
\newcommand{\starh}{h^{*}}
\newcommand{\dd}{\delta}
\DeclareMathOperator{\ord}{ord}
\DeclareMathOperator{\diag}{diag}
\DeclareMathOperator{\Frac}{Frac}
\newtcolorbox{resultbox}{colback=pale,colframe=accent,boxrule=0.6pt,
 arc=1mm,left=9pt,right=9pt,top=7pt,bottom=7pt,breakable}

\begin{document}
\begin{center}
{\small\sffamily\color{accent}INTEGER SEQUENCES\quad /\quad GENERATING FUNCTIONS\quad /\quad ELLIPTIC CURVES}\par
\vspace{15pt}
{\LARGE\bfseries\color{ink}Elliptic addition and\par Jacobi continued fractions}\par
\vspace{9pt}
{\large An explicit proof of a conjectured correspondence,\par with indexing and torsion corrections}\par
\vspace{13pt}
{\normalsize Research report prepared with ChatGPT}\par
{\small 19 September 2026}
\end{center}
\vspace{8pt}
\begin{abstract}
Paul Barry's 2023 paper \emph{Integer sequences from elliptic curves}
conjectures that the Jacobi continued-fraction coefficients of a particular
algebraic generating function are given by coordinates of successive
multiples of a point on an elliptic curve. We give an explicit algebraic proof
of a consistently indexed, non-torsion formulation of that assertion.
The main construction is a closed expression for every normalized complete
quotient. A factorization shows that deleting one continued-fraction layer
is exactly translation by the distinguished elliptic-curve point, after a
rational change of coordinates. No inference from finitely many coefficients
is used in the proof.

The corresponding formulas for both coordinates in terms of Hankel and
modified Hankel determinants follow from a self-contained moment argument.
We also derive the associated Somos-4 recurrence and recover the
already-known division-polynomial identity, explicitly crediting the
2026 proof of Liu, Wang, and Zhang. Torsion specializations explain an essential
missing hypothesis: a zero Hankel determinant need not cause all later ones
to vanish. In fact, the generating function is nonrational whenever the
parameter $b$ is nonzero, even in these exceptional cases. Exact computations
on 106 nonsingular curves, sixteen symbolic identities, and four worked
examples accompany the article. The literature audit identifies the precise
conjecture addressed and distinguishes the present proof from previously
established results; priority is not asserted.
\end{abstract}
\vspace{3pt}
\begin{resultbox}
\textbf{Central correspondence.} On
$Y^2+aXY+bY=X^3+cX^2+dX$, write $mP=(X_m,Y_m)$ for $P=(0,0)$.
For the generating function defined in Section~\ref{sec:setup}, the regular
Jacobi coefficients are
\[
 \alpha_0=1,\quad\alpha_1=-1,\quad\beta_1=1,\qquad
 \alpha_m=\frac{bY_m}{X_m}-d-1,\quad
 \beta_m=-b^2X_m\quad(m\ge2).
\]
The proof constructs all the tails, rather than merely guessing these coefficients.
\end{resultbox}
\clearpage
{\small\setlength{\parskip}{0pt}\tableofcontents}
\clearpage

\section{The problem and the literature boundary}\label{sec:status}

\subsection{The published conjecture}
Barry associates an integer-coefficient generating function to a generalized
Weierstrass equation with a specified integral point. In Section~3 of
\cite{Barry2023}, Conjecture~2 concerns its Hankel transform and an elliptic
division sequence; Conjecture~3 expresses point coordinates through Hankel
determinants; Conjecture~4 proposes a Jacobi continued fraction in the reverse
direction. We take the last assertion as the main target, with the coordinate
formulas as its natural consequence.

This is a tractable kind of sequence problem because the generating function
has degree two over a rational-function field. Its continued-fraction tails
can plausibly remain in the same quadratic extension. The difficulty is not
computing a few tails: it is finding a closed family stable under the
continued-fraction transformation and proving that its parameter evolution
agrees with the elliptic group law.

\subsection{What was already known by the research date}
The current literature matters here. Liu, Wang, and Zhang's preprint
\cite{LWZ2026}, submitted in August 2026, explicitly lists Barry's 2023
Conjecture~2 among the assertions it proves. Its Theorems~5.4--5.5 give the
Somos-4 recurrence and the normalized Hankel/division-polynomial identity.
Those are therefore \emph{not} claimed as new results in this report.

The general relationship among continued fractions, elliptic translation,
orthogonal polynomials, and Hankel determinants is older. Hone's work
\cite{Hone2021} provides a broad treatment, including higher genus.
Barry's 2025 paper \cite{Barry2025} develops related lattice-path and
continued-fraction correspondences, and Chang and Liu \cite{CL2026}
treat further issues involving hyperelliptic functions and bilateral Somos
sequences. The present calculation belongs to this established framework;
it is not a discovery of the general connection.

The specific deliverable here is a directly checkable complete-quotient
formula in the four parameters of \cite{Barry2023}, together with a proof
that it yields the intended assertions of Conjectures~3--4. The sources
inspected did not supply this exact proof in this normalization. That is a
bounded literature observation, not a guarantee that no equivalent result
or unpublished proof exists. Thus the report proves a precise version of a
published conjectural statement, but does not certify a first-in-the-literature
resolution of a still-open problem as of September 2026.

\subsection{Two necessary corrections}
There are two different issues, which should not be conflated.

First, the sentence immediately before Conjecture~4 in \cite{Barry2023}
identifies $(x_n,y_n)$ with $n(0,0)$, while its displayed formulas use
$\beta_{n+1}=-b^2x_n$ and
$\alpha_n=by_{n-1}/x_{n-1}-d-1$. Taken literally, the first nontrivial
coefficient would involve $x_1=0$, and $\alpha_2$ would contain $0/0$.
The same paper's coordinate formulas initialize $(x_0,y_0)=(0,0)$.
The consistent interpretation is
\begin{equation}\label{eq:indexing}
 (x_n,y_n)=((n+1)P),\qquad n\ge0,
\end{equation}
not $nP$. We instead write $(X_m,Y_m)=mP$ throughout, reserving the
point at infinity $\OO$ for $0P$. This removes the ambiguity.

Second, a point can be torsion. Then some multiples reach $\OO$ or
$-P=(0,-b)$, and the displayed affine ratios cease to exist. An infinite
\emph{regular} Jacobi fraction consequently requires a nondegeneracy
hypothesis. Infinite order of $P$ supplies exactly what is needed here.
Section~\ref{sec:torsion} explains what remains valid at torsion points.
A typographical indexing contradiction alone would be an uninteresting
``counterexample''; our goal is to prove the substantive intended formula.

\section{Normalization and the main theorem}\label{sec:setup}

Let $\K$ be a field of characteristic zero. Consider the nonsingular curve
\begin{equation}\label{eq:curve}
 E:\quad Y^2+aXY+bY=X^3+cX^2+dX,
 \qquad P=(0,0),\qquad b\ne0.
\end{equation}
The coefficients $a,b,c,d$ lie in $\K$. Write $mP=(X_m,Y_m)$ whenever
$mP\ne\OO$. Introduce
\begin{align}
 \delta&=abd-b^2c+d^2,& \kappa&=b^4,\label{eq:parameters1}\\
 s&=ab+2d+2,& r&=s-1,& L&=\delta+r.\label{eq:parameters2}
\end{align}
In particular,
\[
 L=(d+1)^2+ab(d+1)-b^2c.
\]
The letter $\delta$ does not denote the discriminant of $E$.

Define the quadratic polynomial and quartic polynomial
\begin{equation}\label{eq:quartic}
 M(z)=1+sz+Lz^2,\qquad
 \D(z)=M(z)^2-4\kappa z^3(1+rz).
\end{equation}
There is a unique $W(z)\in\K[[z]]$ with
\begin{equation}\label{eq:W}
 W(z)^2=\D(z),\qquad W(0)=1.
\end{equation}
Comparing the first four coefficients gives
\begin{equation}\label{eq:Wleading}
 W(z)=1+sz+Lz^2-2\kappa z^3+O(z^4).
\end{equation}

The series $q(z)$ and the moment generating function $G(z)$ are
\begin{align}
 q(z)&=\frac{M(z)-W(z)}{2\kappa z^3},\label{eq:qradical}\\
 G(z)&=\frac{1}{1-z-z^2q(z)}
      =\frac{2\kappa z}{W(z)+2\kappa z(1-z)-M(z)}.
      \label{eq:G}
\end{align}
The apparent powers of $z$ in the denominators cancel. In particular,
$q(0)=G(0)=1$. Equivalently, $q$ is the unique solution with constant term
one of
\begin{equation}\label{eq:qequation}
 M(z)q(z)-\kappa z^3q(z)^2=1+rz.
\end{equation}
This compact normalization is also compatible with the quadratic-series
form used in \cite[Section~5.2]{LWZ2026}. Appendix~\ref{app:normalization}
checks it against Barry's expanded coefficients.

\begin{definition}
A regular Jacobi continued fraction is a formal expansion
\begin{equation}\label{eq:J}
 \cfrac{1}{1-\alpha_0z-
 \cfrac{\beta_1z^2}{1-\alpha_1z-
 \cfrac{\beta_2z^2}{1-\alpha_2z-\cdots}}},
 \qquad \beta_j\ne0\quad(j\ge1).
\end{equation}
Its value is the coefficientwise, or $z$-adic, limit of its finite
truncations. This definition involves no analytic convergence assertion.
\end{definition}

\begin{theorem}[Corrected elliptic Jacobi-fraction formula]\label{thm:main}
Assume that $P$ has infinite order on \eqref{eq:curve}. The generating
function \eqref{eq:G} has the regular Jacobi expansion \eqref{eq:J} with
\begin{equation}\label{eq:coeffs}
 \alpha_0=1,\qquad \alpha_1=-1,\qquad \beta_1=1,
\end{equation}
and, for every $m\ge2$,
\begin{equation}\label{eq:ellipticcoeffs}
 \boxed{\quad
 \alpha_m=\frac{bY_m}{X_m}-d-1,\qquad
 \beta_m=-b^2X_m.\quad}
\end{equation}
These formulas are the consistently indexed version of Barry's
Conjecture~4. The expansion is unique among regular Jacobi fractions.
\end{theorem}

The proof is given in Sections~\ref{sec:map}--\ref{sec:proof}. Its core is
an elementary factorization, not a determinant calculation.

\subsection{Why all denominators exist}
If $P$ has infinite order, then $mP\ne\OO$ for every positive integer $m$.
The only points on \eqref{eq:curve} with $X=0$ are
$P=(0,0)$ and $-P=(0,-b)$. If $mP$ were either of these for $m\ge2$,
then $(m-1)P=\OO$ or $(m+1)P=\OO$, a contradiction. Consequently
\begin{equation}\label{eq:nonzeroX}
 X_m\ne0\quad(m\ge2).
\end{equation}
In particular $\beta_m\ne0$. This simple observation is the exact place
where the infinite-order hypothesis enters the infinite construction.

\subsection{Integer coefficients without elliptic divisions}\label{subsec:recurrences}
Write
\[
 q(z)=\sum_{n\ge0}q_nz^n,\qquad
 G(z)=\sum_{n\ge0}\mu_nz^n.
\]
Equation~\eqref{eq:qequation} gives $q_0=1$, $q_1=-1$, and for $n\ge2$,
\begin{equation}\label{eq:qrec}
 q_n=-s q_{n-1}-Lq_{n-2}
      +\kappa\sum_{i=0}^{n-3}q_iq_{n-3-i},
\end{equation}
where a sum with negative upper limit is zero. The outer fraction gives
$\mu_0=1$ and
\begin{equation}\label{eq:murec}
 \mu_n=\mu_{n-1}+\sum_{i=0}^{n-2}q_i\mu_{n-2-i}\qquad(n\ge1).
\end{equation}
Thus all $q_n$ and $\mu_n$ belong to $\Z[a,b,c,d]$. The rational
elliptic coefficients in \eqref{eq:ellipticcoeffs} produce polynomial
moments; this is a consequence of the algebraic series identity, not an
assumption about cancellation.

For example,
\begin{align*}
 (q_0,q_1,q_2,q_3)&=(1,-1,1-\delta,\delta(s+1)-1+\kappa),\\
 (\mu_0,\ldots,\mu_5)&=(1,1,2,2,4-\delta,
                         4+\kappa+(s-1)\delta).
\end{align*}
The recurrences remain well-defined when the Jacobi formula degenerates.

\section{Elliptic addition as a rational map}\label{sec:map}

\subsection{Coordinates adapted to the fraction}
For a point $Q=(X,Y)$ with $X\ne0$, set
\begin{equation}\label{eq:uv}
 u=\frac{bY}{X}-d-1,\qquad v=-b^2X.
\end{equation}
The inverse transformation is
\begin{equation}\label{eq:uvinverse}
 X=-\frac{v}{b^2},\qquad
 Y=-\frac{v(u+d+1)}{b^3}.
\end{equation}

\begin{lemma}[Invariant cubic]\label{lem:invariant}
The point $Q$ lies on $E$ if and only if
\begin{equation}\label{eq:invariant}
 \inv(u,v):=v^2+(u^2+su+L)v-\kappa(u+1)=0.
\end{equation}
Here the equivalence is on the open set $v\ne0$.
\end{lemma}
\begin{proof}
Put $T=u+d+1$. Substitution of \eqref{eq:uvinverse} into the left side
minus the right side of \eqref{eq:curve}, followed by multiplication by
$b^6/v$, gives
\[
 v^2+v(T^2+abT-b^2c)-b^4(T-d).
\]
The identities
\[
 T^2+abT-b^2c=u^2+su+L,\qquad T-d=u+1
\]
turn this expression into \eqref{eq:invariant}.
\end{proof}

\subsection{The translation law}
The line through $P=(0,0)$ and $Q=(X,Y)$ has slope $t=Y/X$.
Substitution of $Y=tX$ into \eqref{eq:curve} shows that the three
intersection abscissas, with multiplicity, satisfy
\[
 X\bigl(X^2+(c-t^2-at)X+(d-bt)\bigr)=0.
\]
The group law negates the third intersection point. Hence, if
$Q+P=(X',Y')$,
\begin{equation}\label{eq:elladdition}
 X'=t^2+at-c-X,\qquad Y'=-(t+a)X'-b.
\end{equation}
This formula includes tangent intersections when multiplicities occur.

\begin{lemma}[Conjugate translation map]\label{lem:map}
Under \eqref{eq:uv}, translation $Q\mapsto Q+P$ becomes
\begin{align}
 v'&=-v-u^2-su-L=-\frac{\kappa(u+1)}{v},\label{eq:vnext}\\
 u'&=\frac{\kappa}{v'}-s-u,\label{eq:unext}
\end{align}
whenever $v$ and $v'$ are nonzero.
\end{lemma}
\begin{proof}
Multiply the first formula in \eqref{eq:elladdition} by $-b^2$ and use
$b t=u+d+1$. This yields the first expression for $v'$ in
\eqref{eq:vnext}. The second follows by dividing
\eqref{eq:invariant} by $v$.

For the ordinate,
\[
 \frac{bY'}{X'}-d-1
 =-b(t+a)-\frac{b^2}{X'}-d-1
 =-u-s+\frac{\kappa}{v'},
\]
which is \eqref{eq:unext}. Invariance of \eqref{eq:invariant} follows
already from the group law; it can also be checked by direct substitution.
\end{proof}

Thus the geometric orbit provides a sequence
\[
 (u_m,v_m)=(\alpha_m,\beta_m),\qquad m\ge2,
\]
satisfying
\begin{align}
 v_{m+1}+v_m&=-(u_m^2+su_m+L),\label{eq:sumv}\\
 u_{m+1}+u_m&=\frac{\kappa}{v_{m+1}}-s.\label{eq:sumu}
\end{align}
The next section shows that these same relations govern the tails of $G$.

\subsection{The starting point is \texorpdfstring{$2P$}{2P}}
Implicit differentiation of \eqref{eq:curve} at $P$ gives tangent slope
$d/b$. Consequently
\begin{equation}\label{eq:2P}
 X_2=\frac{\delta}{b^2},\qquad
 Y_2=-\left(\frac db+a\right)X_2-b.
\end{equation}
Under the infinite-order hypothesis, $\delta\ne0$ by
\eqref{eq:nonzeroX}. Substitution into \eqref{eq:uv} gives
\begin{equation}\label{eq:initialuv}
 v_2=-\delta,\qquad u_2=-r-\frac{\kappa}{\delta}.
\end{equation}
This initialization is another direct check of the required index shift.

\section{Explicit complete quotients}\label{sec:tails}

\begin{definition}
For a pair $(u,v)$ satisfying \eqref{eq:invariant}, with $v\ne0$, define
\begin{equation}\label{eq:F}
 F_{u,v}(z)=
 \frac{1+sz+(L+2v)z^2-W(z)}
 {2vz^2\bigl(1+(\kappa/v-u)z\bigr)}.
\end{equation}
This is a normalized complete quotient, meaning a tail with constant term one.
\end{definition}

By \eqref{eq:Wleading}, its numerator begins $2vz^2+2\kappa z^3$.
The apparent division by $z^2$ is therefore harmless and
\begin{equation}\label{eq:Fleading}
 F_{u,v}(z)=1+uz+O(z^2).
\end{equation}

\begin{lemma}[The factorization]\label{lem:factor}
Let $N_v(z)=1+sz+(L+2v)z^2$ and $\gamma=\kappa/v-u$.
On the invariant cubic,
\begin{equation}\label{eq:factor}
 N_v(z)^2-W(z)^2
 =4vz^2(1+\gamma z)(1+(s+u)z).
\end{equation}
In particular,
\begin{equation}\label{eq:Finverse}
 \frac1{F_{u,v}(z)}=
 \frac{N_v(z)+W(z)}{2(1+(s+u)z)}.
\end{equation}
\end{lemma}
\begin{proof}
Since $N_v=M+2vz^2$, subtraction using \eqref{eq:quartic} gives
\begin{align*}
 N_v^2-W^2
 &=4vz^2M+4v^2z^4+4\kappa z^3(1+rz)\\
 &=4vz^2\left(1+\left(s+\frac{\kappa}{v}\right)z+
            \left(L+v+\frac{\kappa r}{v}\right)z^2\right).
\end{align*}
The two proposed linear factors have the correct linear coefficient.
Their quadratic coefficient is correct precisely when
\[
 \left(\frac{\kappa}{v}-u\right)(s+u)
       =L+v+\frac{\kappa(s-1)}{v}.
\]
After multiplication by $v$, this equality is equivalent to
\eqref{eq:invariant}. This proves \eqref{eq:factor}.
Multiplying numerator and denominator of the reciprocal of
\eqref{eq:F} by $N_v+W$ now proves \eqref{eq:Finverse}.
All cancelled factors are nonzero formal series or explicit powers of $z$.
\end{proof}

\begin{theorem}[One translation deletes one layer]\label{thm:stripping}
Suppose $(u,v)$ satisfies \eqref{eq:invariant}, and $(u,v)$ and $(u',v')$ are related by
\eqref{eq:vnext}--\eqref{eq:unext}, with $v,v'\ne0$. Then
\begin{equation}\label{eq:stripping}
 \boxed{\qquad
 F_{u,v}(z)=\frac1{1-uz-v'z^2F_{u',v'}(z)}.
 \qquad}
\end{equation}
\end{theorem}
\begin{proof}
Equation~\eqref{eq:unext} gives
\[
 \frac{\kappa}{v'}-u'=s+u.
\]
Thus the denominator of $v'z^2F_{u',v'}$ is
$2(1+(s+u)z)$. By \eqref{eq:Finverse}, the numerator of
$1-uz-1/F_{u,v}$ over that same denominator is
\begin{align*}
 &2(1-uz)(1+(s+u)z)-N_v(z)-W(z)\\
 &\hspace{15pt}=1+sz+
       \bigl(-2u(s+u)-L-2v\bigr)z^2-W(z)\\
 &\hspace{15pt}=1+sz+(L+2v')z^2-W(z).
\end{align*}
The last equality is exactly \eqref{eq:vnext}. This is the numerator of
$v'z^2F_{u',v'}$, proving the identity.
\end{proof}

\begin{remark}
The construction is local and algebraic. The proof of
\cref{thm:stripping} needs the invariant cubic, the displayed rational map,
and nonzero denominators; it does not require an independent theorem about
elliptic functions or a prior Somos recurrence. Nonsingularity and infinite
order are used to identify and continue the geometric orbit globally.
\end{remark}

\section{Proof of the continued-fraction theorem}\label{sec:proof}

\subsection{The initial two layers}
By rationalizing \eqref{eq:qradical},
\begin{equation}\label{eq:qinv}
 \frac1{q(z)}=\frac{M(z)+W(z)}{2(1+rz)}.
\end{equation}
For the initial pair \eqref{eq:initialuv}, one has
\[
 \frac{\kappa}{v_2}-u_2=r.
\]
Hence \eqref{eq:F} specializes to
\begin{equation}\label{eq:F2}
 F_{u_2,v_2}(z)=
 \frac{W(z)-1-sz+(\delta-r)z^2}{2\delta z^2(1+rz)}.
\end{equation}
Direct substitution into \eqref{eq:qinv} gives
\begin{equation}\label{eq:qtail}
 q(z)=\frac1{1+z+\delta z^2F_{u_2,v_2}(z)}
     =\frac1{1-\alpha_1z-\beta_2z^2F_{u_2,v_2}(z)}.
\end{equation}
The definition of $G$ supplies the outer layer,
\begin{equation}\label{eq:Gtail}
 G(z)=\frac1{1-\alpha_0z-\beta_1z^2q(z)},
 \qquad \alpha_0=1,\quad\beta_1=1.
\end{equation}
This proves all the special initial values and links the algebraic series
to the first geometric tail.

\subsection{Iteration and formal convergence}
For every $m\ge2$, put $F_m=F_{u_m,v_m}$, where
$(u_m,v_m)$ is obtained from $mP$. Lemma~\ref{lem:map} and
Theorem~\ref{thm:stripping} give
\begin{equation}\label{eq:Fm}
 F_m=\frac1{1-\alpha_mz-\beta_{m+1}z^2F_{m+1}}.
\end{equation}
All tails exist by \eqref{eq:nonzeroX}. Iterating \eqref{eq:Fm} inside
\eqref{eq:qtail} and \eqref{eq:Gtail} gives the proposed fraction.

For completeness, define
\[
 T_{A,B}(H)=\frac1{1-Az-Bz^2H}.
\]
If $H$ and $\widetilde H$ are formal series with constant term one, then
\begin{equation}\label{eq:valuation}
 T_{A,B}(H)-T_{A,B}(\widetilde H)
 =\frac{Bz^2(H-\widetilde H)}
 {(1-Az-Bz^2H)(1-Az-Bz^2\widetilde H)}.
\end{equation}
Each denominator is a unit. Therefore each surrounding layer increases
an error's order by at least two. Replacing $F_m$ by $1$ makes an
$O(z)$ error at that level, hence an $O(z^{2m+1})$ error in $G$.
The finite expressions stabilize coefficientwise, and their limit is
exactly $G$. This proves existence in \cref{thm:main}.

\subsection{Uniqueness}
If $F=1+f_1z+f_2z^2+\cdots$ is a normalized tail, then any regular
representation $F=(1-Az-Bz^2H)^{-1}$ forces
\[
 A=-[z]F^{-1},\qquad B=-[z^2]F^{-1},\qquad
 H=\frac{1-Az-F^{-1}}{Bz^2}.
\]
Thus the coefficient pair and next tail are uniquely determined whenever
$B\ne0$. Induction proves uniqueness of the entire regular fraction.
Together with the preceding subsections, this completes the proof of
\cref{thm:main}.

\section{Hankel determinants and both point coordinates}\label{sec:hankel}

Let $\mu_n=[z^n]G(z)$. Define
\begin{equation}\label{eq:h}
 h_n=\det(\mu_{i+j})_{0\le i,j\le n},\qquad h_{-1}=1.
\end{equation}
The modified determinant $h_n^*$ has the same first $n$ rows as $h_n$,
but its last row is shifted one place:
\begin{equation}\label{eq:hstar}
 h_n^*=\det A^{(n)},\qquad
 A^{(n)}_{ij}=\begin{cases}
 \mu_{i+j},&i<n,\\
 \mu_{n+j+1},&i=n,
 \end{cases}\qquad 0\le i,j\le n.
\end{equation}
Set $h_{-1}^*=0$. In particular $h_0=1$ and $h_0^*=\mu_1=1$.
The star does not mean that every entry in the matrix has been shifted;
that different convention would give a different determinant.

\subsection{The formal moment lemma}
\begin{lemma}[Jacobi coefficients from determinants]\label{lem:hankel}
For a regular fraction \eqref{eq:J}, with moment sequence $\mu_n$,
\begin{align}
 h_n&=\prod_{j=1}^{n}\beta_j^{\,n+1-j},\label{eq:hproduct}\\
 \beta_n&=\frac{h_nh_{n-2}}{h_{n-1}^2}\qquad(n\ge1),\label{eq:betaH}\\
 \alpha_n&=\frac{h_n^*}{h_n}-\frac{h_{n-1}^*}{h_{n-1}}
                 \qquad(n\ge0).\label{eq:alphaH}
\end{align}
For $n=1$, formula \eqref{eq:betaH} uses $h_{-1}=1$.
\end{lemma}
\begin{proof}
We give a formal proof, so positivity and analytic moment measures are
unnecessary. Let $V$ be the vector space with basis $e_0,e_1,\ldots$, and
let $J$ be the tridiagonal operator
\[
 Je_n=e_{n+1}+\alpha_ne_n+\beta_ne_{n-1},
 \qquad \beta_0e_{-1}:=0.
\]
The coefficient of $e_0$ in $J^ke_0$ is the total weight of all length-$k$
paths on the nonnegative integers starting and ending at zero. An up-step
has weight one, a horizontal step at height $n$ has weight $\alpha_n$,
and a down-step from height $n$ has weight $\beta_n$.
Decomposing paths at their first horizontal step or first up-and-down
excursion gives \eqref{eq:J}. Consequently that coefficient is $\mu_k$.

Set $w_0=1$ and $w_n=\beta_1\cdots\beta_n$. Give $V$ the symmetric
bilinear form $\langle e_i,e_j\rangle=w_i\mathbf1_{i=j}$. Then $J$ is
self-adjoint because $w_{n+1}=\beta_{n+1}w_n$. Let
\[
 \mathcal{L}(f)=\langle f(J)e_0,e_0\rangle;
 \qquad \mathcal{L}(t^k)=\mu_k.
\]
Define monic polynomials by $p_0=1$, $p_1=t-\alpha_0$, and
\[
 p_{n+1}=(t-\alpha_n)p_n-\beta_np_{n-1}.
\]
Induction gives $p_n(J)e_0=e_n$. Self-adjointness therefore implies
\[
 \mathcal{L}(p_ip_j)=w_i\mathbf1_{i=j}.
\]
The basis change from $1,t,\ldots,t^n$ to $p_0,\ldots,p_n$ is
unitriangular. The Gram determinant in the monomial basis is $h_n$,
whereas in the polynomial basis it is $\prod_{j=0}^n w_j$.
This proves \eqref{eq:hproduct}, and its adjacent ratios give
\eqref{eq:betaH}.

Finally, the coefficient of $t^n$ in $p_{n+1}$ is
$-\sum_{j=0}^n\alpha_j$ by the three-term recurrence. The determinant
formula for the monic polynomial orthogonal to $1,t,\ldots,t^n$ is
\[
 p_{n+1}(t)=\frac1{h_n}
 \begin{vmatrix}
 \mu_0&\mu_1&\cdots&\mu_{n+1}\\
 \mu_1&\mu_2&\cdots&\mu_{n+2}\\
 \vdots&\vdots&&\vdots\\
 \mu_n&\mu_{n+1}&\cdots&\mu_{2n+1}\\
 1&t&\cdots&t^{n+1}
 \end{vmatrix}.
\]
Indeed it is monic, and applying $\mathcal L(t^j\cdot)$ for $j\le n$
creates two equal rows. Expansion along the last row shows that its
$t^n$ coefficient is $-h_n^*/h_n$: the relevant minor is the transpose
of \eqref{eq:hstar}. Thus $h_n^*/h_n=\sum_{j=0}^n\alpha_j$.
Subtracting the formula at $n-1$ proves \eqref{eq:alphaH}.
\end{proof}

\subsection{The coordinate theorem}
\begin{corollary}[Corrected form of Barry's Conjecture 3]\label{cor:coordinates}
Under the hypotheses of \cref{thm:main}, all $h_n$ are nonzero. For $m\ge2$,
\begin{align}
 X_m&=-\frac1{b^2}\frac{h_{m-2}h_m}{h_{m-1}^2},\label{eq:XH}\\
 Y_m&=-\frac1{b^3}\frac{h_{m-2}h_m}{h_{m-1}^2}
 \left(\frac{h_m^*}{h_m}-\frac{h_{m-1}^*}{h_{m-1}}+d+1\right).
 \label{eq:YH}
\end{align}
Together with $(X_1,Y_1)=(0,0)$, these recover all positive multiples of $P$.
\end{corollary}
\begin{proof}
Nonvanishing follows from \eqref{eq:hproduct}. Combine
\eqref{eq:ellipticcoeffs} with \eqref{eq:betaH} and \eqref{eq:alphaH},
and then solve the first equation in \eqref{eq:uv} for $Y_m$.
\end{proof}

Replacing $m$ by $n+1$ gives the intended formulas in
\cite[Conjecture~3]{Barry2023}. The star determinants are necessary to
recover the ordinate with its sign; the unmodified Hankel ratios alone
recover only the abscissa.

\section{The known Somos and division-polynomial consequences}\label{sec:somos}

This section gives an independent route from the complete-quotient proof
to results already established in \cite[Theorems~5.4--5.5]{LWZ2026}.
It is included to connect the proof to integer sequences and to control
exceptional specializations, not to claim a new Somos-4 theorem.

\subsection{A second-order recurrence for the Jacobi numerators}
\begin{proposition}\label{prop:betasomos}
For the nondegenerate orbit, the numerator coefficients satisfy
\begin{equation}\label{eq:betasomos}
 v_{m-1}v_m^2v_{m+1}=\kappa^2(v_m+\delta),\qquad m\ge3.
\end{equation}
The lower bound $m\ge3$ is essential: the exceptional value $\beta_1=1$
is not a member of the geometric orbit beginning at $2P$.
\end{proposition}
\begin{proof}
Write $u=u_m$, $v=v_m$, and use \eqref{eq:sumu} at $m-1$:
\[
 u_{m-1}=\frac{\kappa}{v}-s-u.
\]
Since $v_{m-1}v=-\kappa(u_{m-1}+1)$,
\[
 v_{m-1}=\frac{\kappa(s+u-1)}{v}-\frac{\kappa^2}{v^2},
 \qquad v_{m+1}=-\frac{\kappa(u+1)}v.
\]
Multiplication gives
\[
 v_{m-1}v^2v_{m+1}
 =\kappa^2\left(-(u+1)(s+u-1)+\frac{\kappa(u+1)}v\right).
\]
By \eqref{eq:invariant}, the last fraction is $v+u^2+su+L$.
The bracket simplifies to $v+L-s+1=v+\delta$.
\end{proof}

\subsection{The bilinear recurrence and its normalization}
Using \eqref{eq:betaH}, the left side of \eqref{eq:betasomos}, with
$m=n-1$, is
\[
 \frac{h_nh_{n-4}}{h_{n-2}^2}.
\]
Therefore, for $n\ge4$,
\begin{equation}\label{eq:hSomos}
 h_nh_{n-4}=b^8h_{n-1}h_{n-3}+b^8\delta h_{n-2}^2.
\end{equation}
Define
\begin{equation}\label{eq:S}
 S_n=b^{-n^2+2n}h_n.
\end{equation}
The differences between the quadratic exponents on the two sides are
six and eight, respectively. Hence
\begin{equation}\label{eq:SSomos}
 S_nS_{n-4}=b^2S_{n-1}S_{n-3}+\delta S_{n-2}^2.
\end{equation}
The first four values are
\begin{equation}\label{eq:Sinit}
 S_0=1,\quad S_1=b,\quad S_2=-\delta,\quad
 S_3=-b\bigl((s-2)\delta+\kappa\bigr).
\end{equation}
Indeed $h_0=h_1=1$, $h_2=v_2=-\delta$, and
$h_3=v_2^2v_3=-\kappa((s-2)\delta+\kappa)$ by
\eqref{eq:initialuv} and \eqref{eq:vnext}.

\subsection{Identification with division polynomials}
Let $\psi_n$ be the standard division polynomials of \eqref{eq:curve},
and put $D_n=\psi_n(P)$. The initial values needed here are
\begin{equation}\label{eq:Dinit}
 D_0=0,\quad D_1=1,\quad D_2=b,\quad D_3=-\delta,\quad
 D_4=-b\bigl((s-2)\delta+\kappa\bigr).
\end{equation}
For example, in the standard Weierstrass invariants,
$b_4=ab+2d=s-2$, $b_6=b^2$, and $b_8=-\delta$.
Thus $D_3=b_8$ and $D_4=b(b_4b_8-b_6^2)$.
The classical division-polynomial addition identity, recorded in
\cite[Equation~(5.5)]{LWZ2026}, yields
\begin{equation}\label{eq:DSomos}
 D_{j+2}D_{j-2}=b^2D_{j+1}D_{j-1}+\delta D_j^2.
\end{equation}

\begin{corollary}[Known identity, recovered here]\label{cor:division}
For every $n\ge0$,
\begin{equation}\label{eq:division}
 \boxed{\qquad b^{-n^2+2n}h_n=\psi_{n+1}(P).\qquad}
\end{equation}
The identity holds for arbitrary characteristic-zero specializations with
$b\ne0$; nonsingularity is required only when interpreting it as a
statement about the elliptic group law.
\end{corollary}
\begin{proof}
First work over the generic field $\Q(a,b,c,d)$. The generic point $P$
is non-torsion: a hypothetical generic identity $NP=\OO$ would force
$\psi_N(P)$ to be the zero polynomial, contrary to its specialization at
$E:Y^2+Y=X^3-X$, whose point $(0,0)$ has infinite order as shown in
Section~\ref{subsec:small}. The standard characterization
$\psi_j(P)=0$ if and only if $jP=\OO$ is used here.

All terms are therefore nonzero in the generic field. Equations
\eqref{eq:SSomos} and \eqref{eq:DSomos}, with the matching initial
values \eqref{eq:Sinit}--\eqref{eq:Dinit}, identify $S_n$ with $D_{n+1}$
by induction. This proves \eqref{eq:division} generically.

Both sides belong to $\Q[a,b,b^{-1},c,d]$: the moments and their
Hankel determinants are polynomials by \eqref{eq:qrec}--\eqref{eq:murec},
and division polynomials are polynomials in the Weierstrass coefficients
and point coordinates. Equality in the fraction field is consequently
an identity in this ring. It specializes at every $b\ne0$, even when
some later $D_j$ or $h_j$ vanish. No division by a specialized zero is
performed in this argument.
\end{proof}

The proof above also explains why the division-free form
\eqref{eq:hSomos} specializes safely, whereas solving it for $h_n$ by
dividing by $h_{n-4}$ need not do so. This distinction is important at
torsion points.

\section{Torsion, vanishing determinants, and nonrationality}\label{sec:torsion}

\subsection{A complete description of the zero locations}
\begin{corollary}\label{cor:torsion}
Assume $E$ is nonsingular and $b\ne0$. If $P$ has finite order $N$, then
$N\ge3$ and
\begin{equation}\label{eq:zeros}
 h_n=0\quad\Longleftrightarrow\quad N\mid n+1.
\end{equation}
Thus the first zero is $h_{N-1}$, but the next determinant $h_N$ is nonzero.
If $P$ has infinite order, no $h_n$ vanishes.
\end{corollary}
\begin{proof}
Since $-P=(0,-b)\ne P$, the point $P$ is not of order two, and it is
not $\OO$. Apply \eqref{eq:division} and the vanishing characterization
of division polynomials.
\end{proof}

At order $N$, the point $(N-1)P=-P$ has abscissa zero. Formally the
numerator value $-b^2X_{N-1}$ is therefore zero, but the ordinate ratio
$bY_{N-1}/X_{N-1}$ is undefined. There is no valid ordinary Jacobi tail
with these specialized scalar coefficients. It would be incorrect to
interpret the zero numerator by itself as a terminating fraction.

\subsection{Why the fraction cannot terminate}\label{subsec:nonrational}
\begin{proposition}\label{prop:nonrational}
For every characteristic-zero parameter choice with $b\ne0$, the series
$q(z)$ and $G(z)$ in \eqref{eq:qradical}--\eqref{eq:G} are nonrational
over $\K(z)$. This conclusion does not require nonsingularity or
infinite order.
\end{proposition}
\begin{proof}
If $W(z)$ were rational, the equality $W^2=\D\in\K[z]$ would force
$W\in\K[z]$: in a reduced fraction, a nonconstant denominator cannot
have its square divide the square of the numerator. Since $\deg\D\le4$,
one would then have $\deg W\le2$.

But \eqref{eq:Wleading} has cubic coefficient $-2\kappa=-2b^4\ne0$.
This is impossible. Thus $W$ is not rational. Equation~\eqref{eq:qradical}
shows that rationality of $q$ is equivalent to rationality of $W$.
Equation~\eqref{eq:G} shows that rationality of $G$ would imply that of
$q=(1-z-G^{-1})/z^2$. Both series are therefore nonrational.
\end{proof}

A finite ordinary Jacobi fraction is rational. Hence
\cref{prop:nonrational} gives a second, independent reason why torsion
specializations cannot be handled by naive termination. Equivalently, for
$b\ne0$ the moment sequence cannot satisfy any eventual homogeneous
linear recurrence with constant coefficients in $\K$: multiplying a
rational generating function by its denominator gives such a recurrence,
and conversely such a recurrence makes the generating function rational.

\subsection{Two explicit torsion examples}
For
\begin{equation}\label{eq:order3}
 E_3:\quad Y^2+Y=X^3,
\end{equation}
we have $(a,b,c,d)=(0,1,0,0)$, $\delta=0$, and curve discriminant
$-27\ne0$. The tangent at $P$ is $Y=0$, which meets the cubic three
times at $P$. Thus $3P=\OO$. The moment and Hankel sequences begin
\begin{align*}
 (\mu_n)&=(1,1,2,2,4,5,6,15,6,35,25,17,185,\ldots),\\
 (h_n)&=(1,1,0,-1,-1,0,1,1,0,\ldots).
\end{align*}
Here $h_2=0$ but $h_3=-1$. If one set $\beta_2=0$ and truncated after
$\alpha_1=-1$, one would obtain
\[
 \frac1{1-z-z^2/(1+z)}=\frac{1+z}{1-2z^2},
\]
whose coefficient of $z^5$ is $4$, while the actual coefficient is
$\mu_5=5$. The failed specialization is visible at a very small index.

For
\begin{equation}\label{eq:order4}
 E_4:\quad Y^2+XY+Y=X^3+X^2,
\end{equation}
we have $2P=(-1,0)$, which is a nonidentity point of order two. Thus $P$
has order four. The curve discriminant is $-15$ and
\[
 (h_n)=(1,1,1,0,-1,-1,-1,0,1,\ldots).
\]
Again the isolated zero does not propagate to every later determinant.
These examples refute the unqualified extension of the ordinary
nondegenerate fraction to all nonsingular curves; they do not refute
\cref{thm:main}.

\section{Worked non-torsion examples}\label{sec:examples}

\subsection{A small example and a certificate of infinite order}\label{subsec:small}
Take
\begin{equation}\label{eq:small}
 E:\quad Y^2+Y=X^3-X,\qquad P=(0,0).
\end{equation}
Then $\delta=\kappa=1$, $s=0$, $r=-1$, $L=0$, and
\begin{equation}\label{eq:smallG}
 W(z)=\sqrt{1-4z^3+4z^4},\qquad
 q(z)=\frac{1-W(z)}{2z^3},\qquad
 G(z)=\frac{2z}{W(z)+2z-2z^2-1}.
\end{equation}
The discriminant is $37$, so the curve is nonsingular.

There are elementary certificates that $P$ has infinite order.
For a direct certificate using computed point coordinates, the change
of variables
\[
 U=4X,\qquad V=8Y+4
\]
transforms \eqref{eq:small} into $V^2=U^3-16U+16$.
The first few additions give $7P=(-5/9,8/27)$, so its transformed
abscissa is $-20/9$. The Nagell--Lutz theorem says that rational torsion
points on an integral short Weierstrass equation have integral
coordinates; see \cite[Theorem~24.21]{Sutherland2013}. Therefore $7P$,
and hence $P$, has infinite order. This uses a theorem with explicit
hypotheses, not an inference from a long list of nonzero multiples.

The same conclusion can be checked through good reduction:
$\#E(\mathbb F_3)=7$ and $\#E(\mathbb F_5)=8$.
The rational torsion subgroup injects into these good reductions, so its
order divides both numbers and is one; see
\cite[Corollary~24.19]{Sutherland2013}. The verification script counts
these finite sets directly.

\begin{table}[ht]
\centering
\renewcommand{\arraystretch}{1.16}
\begin{tabular}{@{}rrrrr@{}}
\toprule
$m$ & $X_m$ & $Y_m$ & $\alpha_m$ & $\beta_m$\\
\midrule
2 & $1$ & $0$ & $0$ & $-1$\\
3 & $-1$ & $-1$ & $1$ & $1$\\
4 & $2$ & $-3$ & $-3/2$ & $-2$\\
5 & $1/4$ & $-5/8$ & $-5/2$ & $-1/4$\\
6 & $6$ & $14$ & $7/3$ & $-6$\\
7 & $-5/9$ & $8/27$ & $-8/15$ & $5/9$\\
8 & $21/25$ & $-69/125$ & $-23/35$ & $-21/25$\\
\bottomrule
\end{tabular}
\caption{Exact elliptic coordinates and continued-fraction coefficients for
$Y^2+Y=X^3-X$. The special outer values are
$\alpha_0=1$, $\alpha_1=-1$, and $\beta_1=1$.}
\label{tab:small}
\end{table}

The moments and determinants begin
\begin{align*}
 (\mu_n)&=(1,1,2,2,3,4,4,6,7,6,11,10,6,\ldots),\\
 (h_n)&=(1,1,-1,1,2,-1,-3,-5,7,-4,-23,\ldots),\\
 (h_n^*)&=(1,0,0,1,-1,3,2,6,-13,-5,45,\ldots).
\end{align*}
For example, the coordinate formulas at $m=4$ give
\[
 X_4=-\frac{h_2h_4}{h_3^2}=2,\qquad
 Y_4=X_4\left(\frac{h_4^*}{h_4}-\frac{h_3^*}{h_3}\right)
     =2\left(-\frac12-1\right)=-3.
\]
Here $d+1=0$, so there is no additional constant in the ordinate bracket.

\subsection{Barry's larger numerical example}
Barry's motivating example \cite[Section~2]{Barry2023} uses
\[
 (a,b,c,d)=(2,5,4,9),\qquad
 E:Y^2+2XY+5Y=X^3+4X^2+9X.
\]
The discriminant is $-38091$, and direct counting gives
$\#E(\mathbb F_5)=8$ and $\#E(\mathbb F_7)=6$. Both reductions are good,
so the rational torsion subgroup has order dividing two. Since
$-P=(0,-5)\ne P$, the point $P$ cannot be torsion.
Our normalization gives
\[
 \delta=71,\quad\kappa=625,\quad s=30,\quad r=29,\quad L=100,
\]
and hence
\begin{equation}\label{eq:barryG}
 G(z)=\frac{1250z}
 {\sqrt{1+60z+1100z^2+3500z^3-62500z^4}
       +1220z-1350z^2-1}.
\end{equation}
The initial elliptic values are
\[
 2P=\left(\frac{71}{25},-\frac{1974}{125}\right),\qquad
 \alpha_2=-\frac{2684}{71},\qquad \beta_2=-71.
\]
These agree with direct stripping of the moment series. The next point is
\[
 3P=\left(\frac{65325}{5041},\frac{14724605}{357911}\right),
\]
which gives
\[
 \alpha_3=\frac{1089691}{185523},\qquad
 \beta_3=-\frac{1633125}{5041}.
\]
The relatively complicated rational coefficients nevertheless generate the
integer moment sequence
\[
 1,1,2,2,-67,2688,-73696,1856194,-44867225,1067747230,\ldots.
\]
Its normalized Hankel sequence is
\[
 1,5,-71,-13065,-1275214,2876558965,11607013366079,\ldots.
\]
The example is included as a normalization check against the original
paper; all the displayed values are also recomputed independently in the
accompanying files. The main proof is not based on this example.

\section{Exact verification and reproducibility}\label{sec:verification}

\subsection{Independent computational routes}
The file \texttt{verify.py} uses Python's standard library only. Its
calculations use arbitrary-precision integers and exact fractions; no
floating-point tests or numerical tolerances occur. The following routes
are deliberately separate.

The moment series is generated from the division-free quadratic
recurrences \eqref{eq:qrec}--\eqref{eq:murec}. Elliptic multiples are computed
with the generalized Weierstrass chord-and-tangent formulas, not by the
rational map \eqref{eq:vnext}--\eqref{eq:unext}. Jacobi coefficients are
extracted from the moment series by formal reciprocal and coefficient
stripping, not by substituting elliptic coordinates. Hankel and modified
Hankel determinants are computed by fraction-free Bareiss elimination with
row pivoting. Division-polynomial values are computed with the classical
odd-index and doubling recurrences, whose only scalar divisor after
evaluation at $P$ is $D_2=b$. They remain usable when later terms vanish.

The default exhaustive parameter set is
\begin{equation}\label{eq:testgrid}
 a,c,d\in\{-1,0,1\},\qquad b\in\{-2,-1,1,2\}.
\end{equation}
It contains $108$ curves; the two singular ones are explicitly skipped
for the group-law tests. For the remaining $106$ curves, the script
checks all nondegenerate layers up to $m=8$, for a total of $544$
layer checks. It checks \eqref{eq:division} for every $0\le n\le9$
on each of the $106$ curves, including torsion cases.

The four named examples are additionally exported through $\mu_{40}$,
$h_{12}$, $h_{12}^*$, and $12P$. Their division-polynomial identities
are checked through $n=12$. The finite Catalan coefficient formula in
Appendix~\ref{app:coefficients} is independently checked for
$0\le n\le20$ in all four examples, giving $84$ finite-sum checks.
All these tests passed in the delivered run.

\subsection{Symbolic checks and their scope}
The optional file \texttt{verify\_symbolic.py}, tested with SymPy~1.14.0,
verifies sixteen rational-function or polynomial identities. These include
the conversion of Barry's quartic coefficients, the birational invariant
cubic, the complete-quotient factorization, the one-step numerator identity,
the initial tail identity, and the Somos relation for the numerator
coefficients. The program reduces identities on the invariant curve by
substituting
\[
 L=\frac{\kappa(u+1)}v-v-u^2-su,
\]
which is valid on the same open set $v\ne0$ used in the proof.
Every symbolic check passed.

The numerical and symbolic tests are audits of the written argument, not
proof substitutes. They cannot establish novelty, check every future
specialization by enumeration, or replace the nonvanishing and formal-limit
arguments in Sections~\ref{sec:setup} and~\ref{sec:proof}. No Lean or other
proof-assistant formalization is claimed.

\subsection{Running the package}
From the extracted archive's main directory, run
\begin{verbatim}
python verify.py --output data
python verify_symbolic.py       # optional: requires sympy
pdflatex -interaction=nonstopmode -halt-on-error article.tex
pdflatex -interaction=nonstopmode -halt-on-error article.tex
\end{verbatim}
The first command regenerates the exact data and a machine-readable
verification summary. The optional symbolic audit does not affect the
standard-library verification. The two LaTeX runs resolve the contents,
cross-references, and bibliography. A \texttt{Makefile} supplies the same
actions. The bibliography is embedded in \texttt{article.tex}, so a BibTeX
run or network connection is not required.

The quadratic recurrences need $O(N^2)$ arithmetic operations for the first
$N$ moments using direct convolution. Computing all determinants individually
through index $N$ with cubic-time elimination uses $O(N^4)$ arithmetic
operations. These are arithmetic-operation counts, not bit-complexity
claims: integer and rational numerators can become large. The code prioritizes
transparent independent verification rather than a new fastest algorithm.

\section{Conclusion: what has been established}\label{sec:conclusion}

The main result is a proof of the naturally corrected, non-torsion version
of the elliptic Jacobi-fraction formula stated as Conjecture~4 in
\cite{Barry2023}. Its decisive ingredient is the explicit tail
$F_{u,v}$ in \eqref{eq:F}. The invariant equation of the elliptic curve
factors its algebraic norm; that factorization makes continued-fraction
stripping identical to elliptic translation. The argument proves all
coefficients at once and supplies the formal convergence justification.

The Hankel formulas recover both coordinates, giving the correspondingly
indexed form of Conjecture~3. The same calculation recovers the known
Somos-4 and division-polynomial statements, while the universal polynomial
identity explains exactly what survives at torsion specializations. Such
specializations expose a genuine domain issue: ordinary regular Jacobi
fractions break down, but the algebraic generating function remains defined
and nonrational, and Hankel determinants resume being nonzero after an
isolated zero.

The mathematical claim and the historical claim have different scopes.
The complete algebraic proof and its tests are provided here. The broader
elliptic/continued-fraction theory and the normalized Hankel identity are
already part of the literature. Establishing bibliographic priority for
this particular explicit proof would require further specialist review;
it is not inferred merely from the word ``conjecture'' in an older paper.

\appendix
\section{Matching the original expanded generating function}\label{app:normalization}

Barry's generating function is written as
\[
 g(z)=\frac{2b^4z}{\sqrt{Az^4+Bz^3+Cz^2+D_1z+1}+Fz^2+G_1z-1}.
\]
The present parameters give the much shorter coefficient list
\begin{align}
 A&=L^2-4\kappa r,& B&=2sL-4\kappa,\label{eq:AB}\\
 C&=s^2+2L,&D_1&=2s,\label{eq:CD}\\
 F&=-(L+2\kappa),&G_1&=2\kappa-s.\label{eq:FG}
\end{align}
For example,
\begin{align*}
 L^2-4\kappa r
 ={}&a^2b^2(d+1)^2
 -2ab\bigl(2b^4+b^2c(d+1)-(d+1)^3\bigr)\\
 &+b^4\bigl(c^2-4(2d+1)\bigr)
 -2b^2c(d+1)^2+(d+1)^4,
\end{align*}
which is Barry's coefficient $A$. The other identities follow by direct
expansion and are included in the symbolic audit. Thus the proof concerns
the original generating function, not merely another series sharing a
few initial coefficients.

\section{Finite coefficient formulas from Catalan numbers}\label{app:coefficients}

Let
\[
 C(t)=\sum_{k\ge0}C_kt^k=\frac{1-\sqrt{1-4t}}{2t},\qquad
 C_k=\frac1{k+1}\binom{2k}{k}.
\]
Factoring $M$ out of \eqref{eq:qradical} gives the formal identity
\begin{equation}\label{eq:qcat}
 q(z)=\frac{1+rz}{M(z)}
 C\!\left(\frac{\kappa z^3(1+rz)}{M(z)^2}\right)
 =\sum_{k\ge0}C_k\kappa^kz^{3k}
   (1+rz)^{k+1}M(z)^{-2k-1}.
\end{equation}
No convergence condition is needed: the argument of $C$ has order three.
To extract a completely finite expression, put
\begin{equation}\label{eq:T}
 T_{k,m}(s,L)=
 \sum_{\ell=0}^{\lfloor m/2\rfloor}
 (-1)^{m-\ell}
 \binom{2k+m-\ell}{m-\ell}
 \binom{m-\ell}{\ell}s^{m-2\ell}L^\ell.
\end{equation}
Then
\begin{equation}\label{eq:qclosed}
 q_n=\sum_{k=0}^{\lfloor n/3\rfloor}C_k\kappa^k
 \sum_{j=0}^{\min(k+1,n-3k)}
 \binom{k+1}{j}r^j T_{k,n-3k-j}(s,L).
\end{equation}
To verify \eqref{eq:T}, expand $(1+sz+Lz^2)^{-2k-1}$ using the negative
binomial theorem. If $\ell$ factors contribute $Lz^2$, the total number
of nonconstant factors is $m-\ell$, and the remaining $m-2\ell$
contribute $sz$. This gives exactly the displayed summand. Equation
\eqref{eq:qclosed} now follows from \eqref{eq:qcat}.

There is also a finite composition expression for each moment. From
\eqref{eq:G},
\[
 G(z)=\sum_{j\ge0}\frac{z^{2j}q(z)^j}{(1-z)^{j+1}}.
\]
For $n\ge0$, this implies
\begin{equation}\label{eq:muclosed}
 \mu_n=1+
 \sum_{j=1}^{\lfloor n/2\rfloor}
 \sum_{\substack{k_1,\ldots,k_j\ge0\\k_1+\cdots+k_j\le n-2j}}
 \binom{n-j-(k_1+\cdots+k_j)}{j}
 \prod_{i=1}^j q_{k_i}.
\end{equation}
Substituting \eqref{eq:qclosed} removes the auxiliary sequence entirely.
These formulas provide recurrence-free finite sums for the moments, but
\eqref{eq:qrec}--\eqref{eq:murec} are usually more convenient for computation.

\section{Verification details and boundary conventions}\label{app:verification}

The standard-library verifier uses the following division-polynomial
recurrences after evaluation at $P$:
\begin{align*}
 D_{2m+1}&=D_{m+2}D_m^3-D_{m-1}D_{m+1}^3\qquad(m\ge2),\\
 bD_{2m}&=D_m\bigl(D_{m-1}^2D_{m+2}-D_{m-2}D_{m+1}^2\bigr)
 \qquad(m\ge3),
\end{align*}
with the base values \eqref{eq:Dinit}. These are distinct from the
Somos recurrence used in the proof and do not divide by a possibly zero
$D_{m-2}$. Consequently they provide a useful independent test in the
torsion cases.

The curve discriminant used to exclude singular inputs is
\[
 \Delta_E=-b_2^2b_8-8b_4^3-27b_6^2+9b_2b_4b_6,
\]
where
\[
 b_2=a^2+4c,\qquad b_4=ab+2d,\qquad b_6=b^2,\qquad b_8=-\delta.
\]
There is no extra factor of $1728$ in this formula. The parameters
$\Delta_E$ and $\delta$ are kept separate in the code.

In the point-data CSV files, the row $m=0$ is explicitly the point at
infinity. Rows with $X_m=0$ mark the coordinate ratio as undefined; the
special outer Jacobi values $\alpha_0$, $\alpha_1$, and $\beta_1$ are
not inferred from those rows. A stopping point in series stripping is
reported only as a failure of regular scalar stripping at that depth,
not as a claim that the full generating function terminates.

The archive includes the source, compiled article, verification programs,
logs, JSON summaries, and exact CSV tables. It does not include external
paper PDFs, installed fonts, or generated LaTeX auxiliary files.

\begin{thebibliography}{9}
\bibitem{Barry2023}
P. Barry,
\emph{Integer sequences from elliptic curves},
arXiv:2306.05025v2, 3 November 2023.
Sections 2--3, especially Conjectures 2--4.
\url{https://arxiv.org/abs/2306.05025v2}.

\bibitem{LWZ2026}
F. Liu, Y. Wang, and Z. Zhang,
\emph{Hankel Transform and $(\alpha,\beta)$ Somos-4 Sequences},
arXiv:2608.08703v1, submitted 9 August 2026.
Section 5.2, Theorems 5.4--5.5 and Equation (5.5).
\url{https://arxiv.org/abs/2608.08703v1}.

\bibitem{Hone2021}
A. N. W. Hone,
\emph{Continued fractions and Hankel determinants from hyperelliptic curves},
Communications on Pure and Applied Mathematics \textbf{74} (2021),
2310--2347. DOI: 10.1002/cpa.21923.
Preprint arXiv:1907.05204v3.
\url{https://arxiv.org/abs/1907.05204v3}.

\bibitem{Barry2025}
P. Barry,
\emph{Elliptic Curves, Riordan arrays and Lattice Paths},
arXiv:2507.16765v1, 22 July 2025.
\url{https://arxiv.org/abs/2507.16765v1}.

\bibitem{CL2026}
X.-K. Chang and J. Liu,
\emph{Hankel determinants from quadratic orthogonal pairs for hyperelliptic
functions and their applications},
arXiv:2603.11670v1, 2026.
\url{https://arxiv.org/abs/2603.11670v1}.

\bibitem{Sutherland2013}
A. V. Sutherland,
\emph{18.782 Introduction to Arithmetic Geometry, Lecture 24},
Massachusetts Institute of Technology, 3 December 2013.
Corollary 24.19 and Theorem 24.21.
\url{https://math.mit.edu/classes/18.782/2013fa/LectureNotes24.pdf}.
\end{thebibliography}
\end{document}
